\scanpage{20}
Ceci dit, on a des isom. can.
\[
\left\{\begin{aligned}
\unsure{\Phi|\mathcal L}_A&=({}_pA^*)_{\text{ét}}^\vee\\
\unsure{\Phi|\mathcal L}_{A^*}&=({}_pA)_{\text{ét}}^\vee
\end{aligned}\right.
\]
\uline{NB.} Soit $\varphi$ une polarisation de $A$ (plus gén. une \deleted{\unsure{classe de}} section de $\underline{\operatorname{NS}}_{A/S}$). Alors $\varphi$ définit une \deleted{forme alternée sur}\inserted{section \deleted{\unsure{élément de}} $\varphi'$ de} $\bigwedge^2\operatorname{DR}^1(A/S)$, \deleted{\unsure{ou} alors non dég. Cette forme admet les composantes $\underline H^{10}=\underline t_A^\vee$ et $H^{01}=\underline t_{A^*}$ comme sous-espaces isotropes, et définit donc un accouplement
\diagram{\begin{tikzcd}\underline t_A^\vee\otimes\underline t_{A^*}\arrow[r] & \underline{\mathrm O}_S\end{tikzcd}}
i.e. un hom
\diagram{\begin{tikzcd}\underline t_{A^*}\arrow[r] & \underline t_A\end{tikzcd}}.
Cet hom n'est autre que l'hom tangent à
\diagram{\begin{tikzcd}\widetilde\varphi:A\arrow[r,dash] & A^*\end{tikzcd}}}
\[
\bigwedge^{\unsure{2}}\operatorname{DR}^1(A/S)\simeq
\bigwedge^2\underline t_A^\vee+
\overbrace{\underline t_A^{\unsure{\vee}}\otimes\underline t_{A^*}}^{\operatorname{Hom}(\underline t_A,\underline t_{A^*})}
+\bigwedge^{\unsure{2}}\underline t_{A^*}
\]
Ceci dit, $\varphi'$ est dans le terme médian, et comme tel n'est autre que l'hom tangent
\diagram{\begin{tikzcd}\widetilde\varphi':\underline t_A\arrow[r,dash] & \underline t_{A^*}\end{tikzcd}}
de
\diagram{\begin{tikzcd}\widetilde\varphi:A\arrow[r] & A^*\end{tikzcd}}
défini par $\varphi$, qui est déterminé \deleted{\unsure{globalement}}, par \unsure{tensorisation} \deleted{par} $\otimes_{\mathbb F_p}\underline{\mathrm O}_S$, par
\diagram{\begin{tikzcd}[column sep=large]
\underline{\unsure{\Phi|\mathcal L}}_A\arrow[r,dash]\arrow[d,equal] & \underline{\unsure{\Phi|\mathcal L}}_{A^*}\arrow[d,equal]\\
({}_pA^*)_{\text{ét}}^\vee\arrow[r,dash] & ({}_pA)_{\text{ét}}^\vee
\end{tikzcd}}
transposé de
\diagram{\begin{tikzcd}({}_pA)_{\text{ét}}\arrow[r,dash] & ({}_pA^*)_{\text{ét}}\end{tikzcd}}
induit par
\diagram{\begin{tikzcd}{}_p\widetilde\varphi:{}_pA\arrow[r,dash] & {}_pA^*\end{tikzcd}}
$\ldots\ldots$
\mspage{19}
